\chapter{The Bank Markets Division}\label{ch:fm:the-bank-markets-division}

A large US bank's annual report for 2025 allocates \$149.5 billion of its common equity to its commercial and investment bank and reports that
the business earned 18\% on it. The same report explains how the number was chosen: the allocation ``incorporates Basel III Standardized
risk-weighted assets and the global systemically important banks surcharge \dots\ as well as a simulation of capital depletion in a severe stress
environment''. Every desk inside that business answers for its share of the equity, and its share depends on which of those rules binds for it. A
repo desk that looks nearly free by one rule is the most expensive desk in the building by another.

\section{What a markets division sells: flow, financing and structuring}

A bank's markets division (One Quant Book 1, chapter 2) earns its revenue in three ways. \emph{Flow} desks make markets for clients in rates, FX,
credit, equities and commodities and earn the spread and the franchise (Book 1, chapter 2; Book 9, chapter 24). \emph{Financing} desks lend cash
and securities to clients, in repo, securities lending and prime services, and earn a spread on the balance sheet they lend. \emph{Structuring}
desks sell tailored derivatives and structured products and earn a margin that pays for the risk they warehouse (One Quant Book 5, chapter 19).

The three lines use the bank's resources in very different proportions. A flow desk turns its inventory over and uses risk capital for the
market risk it holds; a financing desk holds large, low-risk assets on the balance sheet for a thin spread; a structuring desk holds long-dated,
complex risk that is capital-intensive per dollar of notional. What the division sells, therefore, is priced not in revenue but in revenue per
unit of the constraint each business consumes.

\begin{omfigure}
\begin{tikzpicture}[font=\footnotesize,>=stealth,box/.style={draw,rounded corners,minimum height=0.62cm,minimum width=2.5cm,align=center,
  fill=omDef!10},res/.style={draw,rounded corners,minimum height=0.62cm,minimum width=2.6cm,align=center,fill=omProp!12}]
\node[box] (flow) at (0,2.6) {flow desks\\{\scriptsize spread, franchise}};
\node[box] (fin) at (0,1.3) {financing desks\\{\scriptsize repo, prime}};
\node[box] (str) at (0,0) {structuring\\{\scriptsize tailored risk}};
\node[res] (rwa) at (5.4,2.6) {risk-weighted assets};
\node[res] (le) at (5.4,1.3) {leverage exposure};
\node[res] (st) at (5.4,0) {stress loss};
\draw[->,thick] (flow) -- (rwa); \draw[->,thin,gray] (flow) -- (le);
\draw[->,thick] (fin) -- (le); \draw[->,thin,gray] (fin) -- (rwa);
\draw[->,thick] (str) -- (st); \draw[->,thick] (str) -- (rwa);
\node[draw,rounded corners,fill=omThm!10,minimum width=2.4cm,minimum height=0.62cm,align=center] (eq) at (9.6,1.3) {group equity\\{\scriptsize the largest need}};
\draw[->] (rwa) -- (eq); \draw[->] (le) -- (eq); \draw[->] (st) -- (eq);
\end{tikzpicture}

\omcaption{What each line of a markets division consumes. Flow and structuring desks use risk-weighted assets and stress capacity; financing desks
use balance sheet, measured by leverage exposure. The group holds equity against the largest of its needs; thick arrows mark each line's main
consumption.}\label{fig:fm:the-bank-markets-division:lines}
\end{omfigure}

\section{The capital stack: common equity, leverage and stress buffers}

\begin{definition}[Common equity tier 1 capital]\label{def:fm:the-bank-markets-division:cet1}
A bank's \emph{common equity tier 1 capital}\index{common equity tier 1 capital} (CET1) is its ordinary shareholders' equity, retained
earnings and reserves, less deductions (goodwill and other intangibles, certain deferred tax assets); it is the capital that absorbs losses first
while the bank is a going concern. Its CET1 ratio is CET1 capital over risk-weighted assets.
\end{definition}

\begin{definition}[Leverage ratio, leverage exposure measure]\label{def:fm:the-bank-markets-division:lr}
A bank's \emph{leverage ratio}\index{leverage ratio} is its Tier 1 capital divided by its \emph{leverage exposure measure}\index{leverage
exposure measure}: its on-balance-sheet assets, less amounts deducted from Tier 1 capital, plus certain off-balance-sheet exposures (derivative
potential exposure, securities financing, commitments), with no weighting by risk.
\end{definition}

\begin{definition}[Stress capital buffer, G-SIB surcharge]\label{def:fm:the-bank-markets-division:buffers}
A \emph{stress capital buffer}\index{stress capital buffer} is a capital buffer set from the capital a bank would lose in the supervisor's severe
stress scenario, added to its minimum ratio. A \emph{G-SIB surcharge}\index{G-SIB surcharge} is an additional capital buffer required of a
global systemically important bank, set by a score of its size, interconnectedness, complexity, cross-jurisdictional activity and
substitutability.
\end{definition}

\begin{dated}{2026-09}{One bank's capital stack at the end of 2025}\label{dat:fm:the-bank-markets-division:jpm}
JPMorgan Chase \& Co.\ (Form 10-K for 2025), Standardized approach: CET1 capital \$288.5 billion on RWA of \$1\,981.7 billion, a CET1 ratio of
14.6\% against a requirement of 11.5\% including buffers (a 4.5\% minimum, a \omterm{def:fm:the-bank-markets-division:buffers}{stress capital buffer} of 2.5\%, the floor, in effect through
30 September 2027, and a GSIB surcharge of 4.5\% under method 2; 2.5\% under method 1). Supplementary \omterm{def:fm:the-bank-markets-division:lr}{leverage ratio} 5.8\% on total leverage
exposure of \$5\,302.0 billion, against a 3.0\% minimum plus a 2.0\% buffer. In November 2025 the Federal Reserve, the OCC and the FDIC replaced the
static enhanced-SLR buffers of the largest banks with a buffer of 50\% of the holding company's method 1 GSIB surcharge (capped at 1\% for its bank
subsidiaries). Equity allocated to the commercial and investment bank: \$149.5 billion at the end of 2025, \$166.5 billion from 1 January 2026;
the business's 2025 net income was \$27\,761 million and its ROE 18\%.
\end{dated}

The stack sets two requirements that bind in different places (\cref{fig:fm:the-bank-markets-division:stack}). For the bank of the dated box,
the risk-based rules require Tier 1 capital of 13.0\% of RWA, \$257.6 billion; the leverage rule requires 5.0\% of its leverage exposure, \$265.1
billion. At the level of the whole bank the leverage rule binds, by a small margin; its leverage exposure is 2.7 times its RWA. A markets division
whose financing business makes its leverage exposure four times its RWA faces a leverage rule that binds by a wide one.

\begin{omfigure}
\begin{tikzpicture}
\begin{axis}[width=10.5cm,height=5cm,xbar stacked,bar width=14pt,xmin=0,xmax=16,ytick={0,1},
  yticklabels={{CET1 / risk-weighted assets},{Tier 1 / leverage exposure}},y dir=reverse,ymin=-0.6,ymax=1.6,
  xlabel={\small per cent},tick label style={font=\footnotesize},table/col sep=comma,grid=major,grid style={gray!20},
  legend cell align=left,legend style={font=\scriptsize,at={(0.5,-0.32)},anchor=north,/tikz/every even column/.append style={column sep=8pt}},legend columns=4]
\addplot[fill=omDef!55,draw=omDef] table[y=k,x=minimum]{figdata/desk/05-the-bank-markets-division/stack.csv};
\addplot[fill=omProp!55,draw=omProp] table[y=k,x=buffer]{figdata/desk/05-the-bank-markets-division/stack.csv};
\addplot[fill=omThm!50,draw=omThm] table[y=k,x=gsib]{figdata/desk/05-the-bank-markets-division/stack.csv};
\addlegendimage{only marks,mark=diamond*,mark size=3.5pt,black}
\legend{minimum,{stress or leverage buffer},GSIB surcharge,{actual ratio}}
\end{axis}
\begin{axis}[width=10.5cm,height=5cm,xmin=0,xmax=16,y dir=reverse,ymin=-0.6,ymax=1.6,hide axis,table/col sep=comma]
\addplot[only marks,mark=diamond*,mark size=3.5pt,black] table[y=k,x=actual]{figdata/desk/05-the-bank-markets-division/stack.csv};
\end{axis}
\end{tikzpicture}

\omcaption{One bank's two requirements at the end of 2025: CET1 over risk-weighted assets (a 4.5\% minimum, a 2.5\% \omterm{def:fm:the-bank-markets-division:buffers}{stress capital buffer} and a
4.5\% GSIB surcharge, 11.5\%, against 14.6\% held) and Tier 1 over leverage exposure (3.0\% plus a 2.0\% buffer, against 5.8\%). Data: JPMorgan
Chase \& Co., Form 10-K for 2025.}\label{fig:fm:the-bank-markets-division:stack}
\end{omfigure}

\section{Return on capital by desk}

\begin{definition}[Allocated equity, binding capital constraint]\label{def:fm:the-bank-markets-division:alloc}
A desk's \emph{allocated equity}\index{allocated equity} is the share of the group's equity it is charged with, computed from what it consumes of
each constraint the group faces. The \emph{binding capital constraint}\index{binding capital constraint} of a desk, or of the group, is the
constraint that requires the most equity for it: risk-weighted assets, leverage exposure or stress loss.
\end{definition}

A desk's return on \omterm{def:fm:the-bank-markets-division:alloc}{allocated equity} is its after-tax profit over its \omterm{def:fm:the-bank-markets-division:alloc}{allocated equity}, the desk-level version of the \omterm{def:fm:the-economics-of-a-trading-firm:roe}{return on equity} of
chapter 1, and it is compared with the bank's hurdle rate (One Quant Book 6, chapter 19). The division of the chapter (illustrative) has six desks.
With equity of 13\% of risk-weighted assets, 5\% of leverage exposure, or the desk's stress loss, a tax rate of 25\% and the desks' figures in the
module (\texttt{fm\_bank.DESKS}), their returns are:

\begin{center}
\small
\begin{tabular}{@{}lrrrl@{}}
\toprule
desk & RWA key & leverage key & binding (largest) & binds\\
\midrule
flow rates & 15.4\% & 9.6\% & 9.6\% & leverage\\
repo & 53.8\% & 5.2\% & 5.2\% & leverage\\
prime services & 23.1\% & 13.5\% & 13.5\% & leverage\\
equity derivatives & 18.1\% & 27.5\% & 18.1\% & RWA\\
FX & 28.8\% & 28.1\% & 28.1\% & leverage\\
credit & 10.4\% & 22.5\% & 10.4\% & RWA\\
\bottomrule
\end{tabular}
\end{center}

\begin{omfigure}
\begin{tikzpicture}
\begin{axis}[width=11.5cm,height=5.4cm,ybar,bar width=8pt,ymin=0,ymax=60,ylabel={\small return on allocated equity (\%)},
  xtick={0,1,2,3,4,5},xticklabels={{flow rates},repo,{prime services},{equity derivatives},FX,credit},
  xticklabel style={font=\scriptsize,text width=1.5cm,align=center},tick label style={font=\footnotesize},table/col sep=comma,
  enlarge x limits=0.1,grid=major,grid style={gray!20},legend cell align=left,
  legend style={font=\scriptsize,at={(0.5,-0.32)},anchor=north,/tikz/every even column/.append style={column sep=8pt}},legend columns=3]
\addplot[fill=omDef!55,draw=omDef] table[x=k,y=rwa]{figdata/desk/05-the-bank-markets-division/roae.csv};
\addplot[fill=omThm!50,draw=omThm] table[x=k,y=leverage]{figdata/desk/05-the-bank-markets-division/roae.csv};
\addplot[fill=gray!40,draw=gray] table[x=k,y=binding]{figdata/desk/05-the-bank-markets-division/roae.csv};
\draw[omProp,dashed,thick] (axis cs:-0.6,12) -- (axis cs:5.6,12);
\legend{equity by RWA,equity by leverage exposure,equity by the binding constraint}
\end{axis}
\end{tikzpicture}

\omcaption{The six desks' after-tax returns on equity allocated by risk-weighted assets, by leverage exposure, and by whichever requires more;
the dashed line is a 12\% hurdle rate. Illustrative desks; requirement levels at the scale of the dated box. Data:
\texttt{fm\_bank.roae\_table}.}\label{fig:fm:the-bank-markets-division:roae}
\end{omfigure}

The repo desk is the division's best business by risk-weighted assets and its worst by leverage. Its assets are short-term secured loans with low
risk weights, and they fill the balance sheet: \$400 billion of leverage exposure for \$15 billion of RWA. Equity derivatives and credit are the
opposite: risk-heavy and balance-sheet-light, they are bound by RWA.

\begin{proposition}[Allocating by each desk's own maximum over-allocates]\label{prop:fm:the-bank-markets-division:max}
Let desk $i$ need $a_i$ of equity under one constraint and $b_i$ under another, and let the group hold $E=\max(\sum_ia_i,\sum_ib_i)$. Then
$\sum_i\max(a_i,b_i)\ge E$, with equality only if the same constraint binds for every desk.
\end{proposition}

\begin{proof}
$\max(a_i,b_i)\ge a_i$ and $\ge b_i$ for every $i$; sum each inequality. Equality in both sums requires $\max(a_i,b_i)=a_i$ for all $i$ or
$=b_i$ for all $i$.
\end{proof}

The chapter's desks need \$35.1 billion of equity by RWA and \$55.5 billion by leverage; the group holds \$55.5 billion, but the desks' own
maxima add up to \$62.1 billion. Charged each by its own binding constraint, the division would report a return of 11.4\% on equity it does not
have. A consistent allocation scales one key to the group's binding constraint, or blends the keys with weights that add to the group's equity,
as the bank of the dated box does when it combines RWA, its GSIB surcharge and a stress simulation.

\section{The balance sheet as a scarce input}

\begin{definition}[Balance-sheet charge]\label{def:fm:the-bank-markets-division:charge}
A \emph{balance-sheet charge}\index{balance-sheet charge} is an internal price per unit of balance sheet, usually of leverage exposure, charged
to a desk for what it uses: the cost of the equity the leverage requirement makes the group hold against it.
\end{definition}

\begin{method}[Setting a balance-sheet charge]\label{met:fm:the-bank-markets-division:charge}
\begin{enumerate}
\item Take the equity the leverage requirement holds per unit of leverage exposure (5\% here, or the bank's target including buffers).
\item Multiply by the hurdle rate: at 12\% the charge is $0.12\times0.05=0.6\%$ a year of leverage exposure.
\item Deduct the charge from each desk's pre-tax profit, alongside funds-transfer pricing (One Quant Book 6, chapter 24) for its funding.
\item A desk whose profit after the charge is negative does not earn its cost of capital on the balance sheet it uses.
\end{enumerate}
\end{method}

At 0.6\% a year the repo desk's \$400 billion of leverage exposure costs \$2.4 billion, and its after-tax profit of \$1.05 billion becomes a loss of
\$0.75 billion. It breaks even at a charge of 0.35\% (35 basis points) a year. The prime services desk earns \$0.45 billion after the charge and
equity derivatives \$1.11 billion. The charge does not say that repo should close; it says that every dollar of balance sheet it uses must be
priced at 60 basis points a year, and that the business is viable only at a spread above that.

\omcode{code/firm/bankdesk/firm_bankdesk.py}{45}{79}{Equity by each constraint and by the binding one, returns on it, the balance-sheet charge,
and the desk mix under the group's constraints.}\label{lst:fm:the-bank-markets-division:bankdesk}

With the group's equity fixed at the \$55.5 billion its leverage requirement holds today, the desk mix that maximises the division's after-tax
profit (a linear programme, \texttt{firm.bankdesk.optimal\_mix}, each desk scaled between zero and 1.5 times its size) shrinks repo to about a
tenth of its size and grows every other desk to the cap (\cref{fig:fm:the-bank-markets-division:mix}): profit rises from \$7.05 billion to \$9.12
billion. The programme assumes each desk can grow without lowering its margins and that the desks are independent. The second assumption is the
one a franchise breaks.

\begin{omfigure}
\begin{tikzpicture}
\begin{axis}[width=11.5cm,height=5cm,ybar,bar width=10pt,ymin=0,ymax=1.7,ylabel={\small scale of the desk (today $=1$)},
  xtick={0,1,2,3,4,5},xticklabels={{flow rates},repo,{prime services},{equity derivatives},FX,credit},
  xticklabel style={font=\scriptsize,text width=1.5cm,align=center},tick label style={font=\footnotesize},table/col sep=comma,
  enlarge x limits=0.1,grid=major,grid style={gray!20},legend cell align=left,
  legend style={font=\scriptsize,at={(0.5,-0.32)},anchor=north,/tikz/every even column/.append style={column sep=8pt}},legend columns=2]
\addplot[fill=gray!40,draw=gray] table[x=k,y=today]{figdata/desk/05-the-bank-markets-division/mix.csv};
\addplot[fill=omDef!55,draw=omDef] table[x=k,y=optimal]{figdata/desk/05-the-bank-markets-division/mix.csv};
\legend{today,profit-maximising mix under the group's two constraints}
\end{axis}
\end{tikzpicture}

\omcaption{The desk mix that maximises the division's after-tax profit with today's equity, when each desk may shrink to zero or grow by half and
margins do not change: the balance-sheet-heavy repo desk gives its leverage exposure to the others. Data: \texttt{fm\_bank.mix},
\texttt{firm.bankdesk.optimal\_mix}.}\label{fig:fm:the-bank-markets-division:mix}
\end{omfigure}

\section{The client franchise}

A client who finances its positions with a bank also clears, borrows and trades with it. The flow desks' franchise (One Quant Book 1, chapter 2)
is built partly on the financing desks' balance sheet: prime services clients choose the bank that lends to them, and trade with it; the central
risk book (One Quant Book 9, chapter 25) nets the flows of all of them. A desk that looks value-destroying alone may be what keeps clients in the
building.

\begin{example}[Cutting repo]\label{ex:fm:the-bank-markets-division:cut}
The division halves its repo desk, and prime services loses 20\% of its revenue, with its costs unchanged, as some clients follow their financing
elsewhere. After-tax profit falls from \$7.05 billion to \$6.08 billion; the equity the desks need falls from \$62.1 billion to \$50.1 billion (each
by its binding constraint); the return on that equity rises from 11.4\% to 12.1\%. Whether the cut is right depends on what the freed equity earns:
returned to shareholders, or given to the desks the linear programme wanted to grow.
\end{example}

\begin{remark}[Why banks specialise]\label{rem:fm:the-bank-markets-division:special}
Because each bank's binding constraint differs, the same business earns different returns in different banks: a bank bound by leverage prices repo
dearly and equity derivatives cheaply; a bank bound by risk-weighted assets does the reverse. The rules are the same for all large banks; their
balance sheets are not, and the businesses each bank wins follow its constraints (One Quant Book 17 surveys the result).
\end{remark}

\section{Tutorial: six desks and three keys}\label{tut:fm:the-bank-markets-division:tut}

\begin{tutorial}
\textbf{Goal.} Allocate a division's equity by each constraint, compute each desk's return, price the balance sheet, and find the best mix.
\textbf{End state:} the table of returns and \cref{fig:fm:the-bank-markets-division:roae,fig:fm:the-bank-markets-division:mix}.
\begin{tutsteps}
\item \textbf{The bank.} \texttt{fm\_bank.jpm\_requirements()} reads the published figures and computes the Tier 1 each rule requires:
\$257.6 billion by RWA, \$265.1 billion by leverage exposure.
\item \textbf{Keys.} \texttt{firm.bankdesk.allocate(DESKS, KEYS)} gives each desk's equity by RWA, leverage and stress, and the largest;
\texttt{roae} divides after-tax profit by each (\cref{lst:fm:the-bank-markets-division:bankdesk}).
\item \textbf{The charge.} \texttt{fm\_bank.charged\_profits()} deducts 0.6\% of leverage exposure; \texttt{breakeven\_charge} finds each desk's
break-even.
\item \textbf{The mix.} \texttt{fm\_bank.mix()} solves the linear programme; \texttt{franchise()} runs \cref{ex:fm:the-bank-markets-division:cut}.
\end{tutsteps}
\end{tutorial}

\textbf{What to change next.} Raise the leverage requirement by the new eSLR buffer and redo the allocation; add a constraint that prime services
cannot exceed repo's size, a crude franchise link, and rerun the programme.

\section{Build: the division's capital model}\label{bld:fm:the-bank-markets-division:bankdesk}

\begin{build}
\bfield{Purpose} The capital side of a bank's markets division: what each desk costs in equity under each rule, what it earns on it, and what the
balance sheet is worth.
\bfield{Interface} \texttt{firm.bankdesk}: \texttt{Desk(name, revenue, cost, rwa, le, stress)}; \texttt{Keys(k\_rwa, k\_le, tax)};
\texttt{allocate}, \texttt{roae}, \texttt{binding}; \texttt{balance\_sheet\_charge}, \texttt{breakeven\_charge};
\texttt{optimal\_mix(desks, equity, keys, x\_max)}.
\bfield{Rules} Each requirement is linear in the desk's size; the binding allocation is reported with a warning that it sums above the group's
equity (\cref{prop:fm:the-bank-markets-division:max}); the linear programme's constraints are the group's, never the desks'.
\bfield{Acceptance tests} \texttt{code/firm/bankdesk/tests/}: equity by key and binding constraint on a two-desk example; the break-even charge
zeroes profit; the optimal mix respects both constraints and favours the balance-sheet-light desk.
\bfield{Stretch} A stress constraint in the programme; a blended allocation that adds to the group's equity; FRTB capital from One Quant Book 6's
\texttt{firm.frtb} as the RWA input for the trading desks.
\end{build}

\begin{omsources}
\item JPMorgan Chase \& Co., Form 10-K for 2025 (segment results, line-of-business equity, capital risk management).
\item The Goldman Sachs Group, Inc., Form 10-K for 2025 (segment results).
\item One Quant Book 6, chapters 19, 23 and 24 (hurdle rate, trading-book capital, funds-transfer pricing).
\end{omsources}

\section{Exercises}

\begin{exercise}[$\star$]\label{exo:fm:the-bank-markets-division:1}
The bank of the dated box had Tier 1 capital of \$307.6 billion and total leverage exposure of \$5\,302.0 billion. Compute its supplementary \omterm{def:fm:the-bank-markets-division:lr}{leverage
ratio} and the Tier 1 its 5.0\% requirement holds.
\end{exercise}

\begin{exercise}[$\star$]\label{exo:fm:the-bank-markets-division:2}
Divide the commercial and investment bank's \$27\,761 million of 2025 net income by its \$149.5 billion of \omterm{def:fm:the-bank-markets-division:alloc}{allocated equity} and compare with the
reported ROE of 18\%. Why may they differ?
\end{exercise}

\begin{exercise}[$\star$]\label{exo:fm:the-bank-markets-division:3}
The repo desk earns \$2.0 billion of revenue for \$0.6 billion of costs, uses \$15 billion of RWA and \$400 billion of leverage exposure. Compute its
after-tax \omterm{def:fm:the-economics-of-a-trading-firm:roe}{return on equity} by each key, at 13\% and 5\% and a 25\% tax rate.
\end{exercise}

\begin{exercise}[$\star\star$]\label{exo:fm:the-bank-markets-division:4}
At a 12\% hurdle rate and a 5\% leverage requirement, what is the \omterm{def:fm:the-bank-markets-division:charge}{balance-sheet charge}, and at what charge does the repo desk break even?
\end{exercise}

\begin{exercise}[$\star\star$]\label{exo:fm:the-bank-markets-division:5}
Show that the desks' own-maximum allocations add up to \$62.1 billion while the group holds \$55.5 billion, and explain the gap with
\cref{prop:fm:the-bank-markets-division:max}.
\end{exercise}

\begin{exercise}[$\star\star$]\label{exo:fm:the-bank-markets-division:6}
In \cref{fig:fm:the-bank-markets-division:roae}, which desks clear a 12\% hurdle under the binding allocation, and which only under the RWA key?
\end{exercise}

\begin{exercise}[$\star\star\star$]\label{exo:fm:the-bank-markets-division:7}
\emph{Coding.} Rerun \texttt{fm\_bank.franchise} with prime services losing 40\% of its revenue instead of 20\%. What happens to profit and to the
return on the desks' equity?
\end{exercise}

\begin{exercise}[$\star\star\star$]\label{exo:fm:the-bank-markets-division:8}
\emph{Find the flaw.} ``The linear programme says repo should shrink by 89\%. Let us do it this quarter.''
\end{exercise}

\section{Problem: The Desk That Uses No Capital}

\begin{problem}[{Weekend problem --- the desk that uses no capital}]\label{pb:fm:the-bank-markets-division:1}
The head of a markets division is told that the repo desk ``uses no capital'' and earns 54\% on it. You are asked to check.

\medskip\noindent\textbf{Part I --- The rules.}
\begin{enumerate}
\item Define CET1 capital, the \omterm{def:fm:the-bank-markets-division:lr}{leverage ratio} and the \omterm{def:fm:the-bank-markets-division:lr}{leverage exposure measure}.
\item Decompose the published bank's 11.5\% CET1 requirement.
\item Compute the Tier 1 its risk-based and its leverage requirements hold, and say which binds.
\item What did the November 2025 rule change in the leverage buffers of the largest banks?
\item How does the published bank allocate equity to its lines of business?
\end{enumerate}

\medskip\noindent\textbf{Part II --- The desks.}
\begin{enumerate}[resume]
\item Define \omterm{def:fm:the-bank-markets-division:alloc}{allocated equity} and the \omterm{def:fm:the-bank-markets-division:alloc}{binding capital constraint}.
\item Give the repo desk's return by RWA, by leverage and by the binding constraint.
\item Which desks are bound by RWA, and why?
\item State and prove \cref{prop:fm:the-bank-markets-division:max}.
\item Give the equity by RWA, by leverage and by desk-level maxima, and the division's return on the last.
\end{enumerate}

\medskip\noindent\textbf{Part III --- The balance sheet.}
\begin{enumerate}[resume]
\item Define a \omterm{def:fm:the-bank-markets-division:charge}{balance-sheet charge} and compute it at a 12\% hurdle.
\item Give each desk's after-tax profit after the charge.
\item At what charge does the repo desk break even?
\item What does the linear programme do with today's equity, and what profit does it reach?
\item Which two assumptions of the programme are the weakest?
\end{enumerate}

\medskip\noindent\textbf{Part IV --- The franchise.}
\begin{enumerate}[resume]
\item What happens to profit, equity and return when repo is halved and prime services loses 20\% of its revenue?
\item Why may a bank bound by leverage and one bound by RWA run different businesses?
\item What would you ask the prime services desk before cutting repo?
\item State the \emph{named result}: the repo desk's return under the RWA and the leverage keys, and its break-even \omterm{def:fm:the-bank-markets-division:charge}{balance-sheet charge}.
\item In two sentences, answer the claim that repo ``uses no capital''.
\end{enumerate}
\end{problem}

\section{Interview questions}

\begin{interviewq}[$\star$ \iqroles{bank}]\label{iq:fm:the-bank-markets-division:1}
What is the difference between a risk-based capital ratio and a \omterm{def:fm:the-bank-markets-division:lr}{leverage ratio}, and why do regulators impose both?
\end{interviewq}

\begin{interviewq}[$\star$ \iqroles{bank, trader}]\label{iq:fm:the-bank-markets-division:2}
Why is a repo desk capital-light by risk-weighted assets and heavy by leverage?
\end{interviewq}

\begin{interviewq}[$\star\star$ \iqroles{bank}]\label{iq:fm:the-bank-markets-division:3}
A desk earns \$1 billion after tax and uses \$200 billion of leverage exposure. At a 5\% leverage requirement and a 12\% hurdle, does it create
value?
\end{interviewq}

\begin{interviewq}[$\star\star$ \iqroles{bank, risk}]\label{iq:fm:the-bank-markets-division:4}
Why can allocating each desk equity by its own binding constraint overstate the equity the group needs?
\end{interviewq}

\begin{interviewq}[$\star\star$ \iqroles{bank}]\label{iq:fm:the-bank-markets-division:5}
Two banks face identical rules. Why might one of them run a large repo business and the other a large equity-derivatives business?
\end{interviewq}

\begin{interviewq}[$\star\star\star$ \iqroles{bank, researcher}]\label{iq:fm:the-bank-markets-division:6}
You are asked to set the desk mix for next year with a fixed amount of equity. What would you add to a linear programme before trusting its
answer?
\end{interviewq}
